\begin{lemma}{Integrierbarkeit von Extrema in Gordon'schen Paaren}
             {IntegrierbarkeitVonExtremaInGordonschenPaaren}
Seien $c \in \R$, $(X, Y)$ das Gordon'sche Paar zum Tupel
$(\gamma, \delta, g, c)$.\\
Seien $\emptyset \neq T \subseteq \R^m$ beschränkt
und $\emptyset \neq V \subseteq \R^N$ beschränkt.\\
Dann gilt:
\begin{enumerate}[(1)]
	\item \label{IntegrierbarkeitVonExtremaInGordonschenPaaren1}
	$\Inf{t \in T}\Sup{v \in V}X(t, v)$ und $\Sup{t \in T}\Sup{v \in V}X(t, v)$
	sind integrierbar.\absatz
	Insbesondere sind $\Inf{t \in T}\bprod t \delta_m$
	und $\Sup{t \in T}\bprod t \delta_m$ integrierbar.\\
	Außerdem ist für jeden
	normierten $\R$--Vektorraum $(E, \norm\cdot_E)$ und $x \in E^N$
die Abbildung $\norm\cdot_E \circ \bprod \cdot x_n \circ \gamma$ integrierbar.
	\item \label{IntegrierbarkeitVonExtremaInGordonschenPaaren2}
	$\Inf{t \in T}\Sup{v \in V}Y(t, v)$ und $\Sup{t \in T}\Sup{v \in V}Y(t, v)$
	sind integrierbar.\\
	Sind $(E, \norm\cdot_E)$ ein normierter $\R$--Vektorraum, 
	$x \in E^N$, $G_x$ der Gauß'sche Operator zum Tupel $(x,g)$, so sind
	\[\inf_{t \in T}\, \norm\cdot_E \circ ev(t,\cdot)\circ G_x \quad\text{ und } 
	  \quad \sup_{t \in T}\, \norm\cdot_E \circ ev(t,\cdot)\circ G_x \quad
	\text{integrierbar.}\]
\end{enumerate}
\end{lemma}
\begin{proof}
Wir zeigen, dass für alle $Z \in \meng{X, Y}$ ein
 $f_Z: \Omega \rightarrow \R$ integrierbar existiert mit
\[(\ast)\quad
\forall\, t \in T \, \forall\, v \in V \, \forall\, \omega \in \Omega:
 \abs{Z(t, v)(\omega)} \le f_Z(\omega)\text.\]
Ist nämlich $Z \in \meng{X, Y}$ und $f_Z :\Omega \rightarrow \R$ integrierbar, 
sodass $(\ast)$ erfüllt ist, so gilt für alle $\omega \in \Omega$:
$\forall\,t \in T \, \forall\, v \in V:
-f_Z(\omega) \le Z(t, v) (\omega) \le f_Z(\omega)$\\
und damit sogleich
\[-f_Z(\omega) \le \inf_{t \in T}\inf_{v \in V} Z(t, v)(\omega)
\le \inf_{t \in T}\sup_{v \in V} Z (t, v) (\omega)
\le \sup_{t \in T}\sup_{v\in V} Z(t, v)(\omega) \le f_Z(\omega)\text.\]
Somit gilt insbesondere
\[-f_Z \le \inf_{t \in T}\sup_{v \in V} Z (t, v)
\le \sup_{t \in T}\sup_{v\in V} Z(t, v) \le f_Z\text.\]
Da $f_Z$ und $-f_Z$ integrierbar und nach
\ref{MessbarkeitVonExtremaInGordonschenPaaren} die Abbildungen
\[\inf_{t\in T}\sup_{v\in V}Z(t, v),\quad \sup_{t \in T}\sup_{v\in V} Z(t, v)\]
messbar sind, sind die letztgenannten Abbildungen nach
\refclaim{Integrierbarkeitskriterien}{Integrierbarkeitskriterien2}
integrierbar.\\
Da $T$ beschränkt ist, gibt es ein $a \in \R_{>0}$ mit 
$\forall\, t \in T: \norm t_{2, m} \le a$.\\
Da auch $V$ beschränkt ist, gibt
es ein $b \in \R_{>0}$ mit $\forall\, v \in V: \norm v_{2, N} \le b$.\\
(1) Sei 
$f_X:=a\abs c(\norm\cdot_{2,m}\circ\delta) + \norm\cdot_{2, N} \circ\gamma$.\\
Nach \refclaim{ElementareEigenschaftenQuadratintegrierbarerAbbildungen}
{ElementareEigenschaftenQuadratintegrierbarerAbbildungen4} sind
$\norm\cdot_{2,m} \circ \delta$ und $\norm\cdot_{2, N} \circ\gamma$ 
$\mfa-\mfb$--messbar und nach
\ref{ErwartungswertDer2NormEinesStandardnormalverteiltenProzesses} sind diese
beiden Abbildungen auch integrierbar.\\
Damit ist es auch $f_X$ als Linearkombination integrierbarer Abbildungen.\\
Seien nun $t \in T$, $v \in V$ und $\omega \in \Omega$.\\
Dann gilt:
\begin{align*}
\abs{X(t, v)(\omega)}&\overset{\refclaim{BeschraenktheitGordonscherPaare}
{BeschraenktheitGordonscherPaare1}}\le
\abs c \norm t _{2, m}\norm {\delta(\omega)}_{2, m} 
+ \norm v _{2, N}\norm{\gamma(\omega)}_{2, N}\\
&\le\abs ca\norm{\delta(\omega)}_{2,m}+b\norm{\gamma(\omega)}_{2,N}=f_X(\omega)
\text.\end{align*}
Damit ist (1) (ohne den Zusatz) nach obiger Begründung gezeigt.\\
Sei $(X', Y')$ das Gordon'sche Paar zum Tupel $(\gamma, \delta, g, 1)$ und
$V' := \meng 0$.\\
Dann ist nach (1) die Abbildung
\[\sup_{t \in T}\bprod t \delta_m =
\sup_{t \in T}\bprod t \delta_m + \sup_{v \in V'} \bprod v \gamma_N
\overset{\refclaim{BeschraenktheitGordonscherPaare}
{BeschraenktheitGordonscherPaare4}}=
\sup_{t \in T}\sup_{v\in V'}X'(t, v) \text{ integrierbar.}\]
Analog ist
\[\inf_{t \in T}\bprod t \delta _m =\inf_{t \in T}\bprod t \delta_m + \sup_{v \in V'} \bprod v \gamma_N
\overset{\refclaim{BeschraenktheitGordonscherPaare}
{BeschraenktheitGordonscherPaare4}}=
\inf_{t \in T}\sup_{v\in V'}X'(t, v) \text{ integrierbar.}\]
Seien nun $(E, \norm\cdot_E)$ ein normierter $\R$-Vektorraum, $x \in E^N$,
$T':= \meng 0$, sowie $V'':= \menge{\pfeil \vi N (x)}{\vi \in \kabg{E'}01}$.\\
Wegen
\[\forall\, \vi \in \kabg{E'}01: \abs{\pfeil\vi N (x)}
\overset{\ref{EigenschaftenDerErweiterungImNormiertenFall}}\le
\opnorm \vi E \R \norm x _{1, N, E} \le \norm x _{1, N, E}\]
ist $V''$ beschränkt und offensichtlich nichtleer.\\
Also gilt für alle $\omega \in \Omega$
\begin{align*}
&\,\left(\norm\cdot_E \circ \bprod \cdot x \circ \gamma \right)(\omega)
=\norm{\bprod {\gamma(\omega)} x_N}_E
 \overset{\ref{NormUndDualeNorm}}= 
\sup_{\vi \in \kabg{E'}01}\vi\left(\bprod{\gamma(\omega)}x_N \right)\\
\overset{\refclaim{EigenschaftenDesLinearkombinators}
{EigenschaftenDesLinearkombinators2}}= &\,
\sup_{\vi \in \kabg{E'}01}\sprod{\gamma(\omega)}{\pfeil \vi N (x)}_N
= \sup_{v \in V''}\sprod{\gamma(\omega)}v_N
=  \sup_{v \in V''}\sprod v{\gamma(\omega)}_N\\
=&\, \left(\sup_{v \in V''}\bprod v \gamma_N\right)(\omega)
\end{align*}
und daher
\begin{align*}
\norm\cdot_E \circ \bprod \cdot x \circ \gamma &=
\sup_{v \in V''}\bprod v \gamma_N 
= \sup_{t \in T'}\bprod t \delta_m + \sup_{v \in V''}\bprod v \gamma_N\\
&\overset{\refclaim{BeschraenktheitGordonscherPaare}
{BeschraenktheitGordonscherPaare4}}= \sup_{t \in T'}\sup_{v \in V''}X'(t,v)
\ \text{ integrierbar nach (1).}
\end{align*}
(2) Die Abbildung $\Mat_{m, N}$ (vgl. Bezeichnungen.(\ref{Bez12})) ist linear
und bijektiv.\\
Da die $\R$--Vektorräume $\R^{mN}$ und $\R^{m\times N}$ endlichdimensional
sind, ist $\Mat_{m,N}$ als Abbildung von $(\R^{mN}, \norm\cdot_{2, mN})$ nach
$(\R^{m \times N}, \opnorm \cdot {\R^m}{\R^N})$ stetig.\\
Wegen der Injektivität von $\Mat_{m,N}$ ist nach
\refclaim{EigenschaftenDerBildhalbnorm}{EigenschaftenDerBildhalbnorm2} 
\[\norm\cdot_{\Mat_{m, N}}: \R^{mN} \rightarrow [0, \infty[,
x \mapsto \opnorm {\Mat_{m, N}(x)}{\R^m}{\R^N}\]
eine Norm auf dem $\R^{mN}$ und da $\Mat_{m, N}$ stetig ist, existiert nach
\refclaim{EigenschaftenDerBildhalbnorm}{EigenschaftenDerBildhalbnorm3} 
ein $d \in \R_{>0}$ mit
$\norm\cdot_{\Mat_{m, N}} \le d \norm\cdot_{2, mN}$.\hfill $(\ast\ast)$\\
Sei nun $f_Y := abd \left(\norm\cdot_{2, mN} \circ \widetilde g \right)$.\\
Da $\widetilde g$ nach Voraussetzung $N(0, 1)$--verteilt und unabhängig ist,
ist $\norm\cdot_{2, mN} \circ \widetilde g$ nach
\refclaim{ElementareEigenschaftenQuadratintegrierbarerAbbildungen}
{ElementareEigenschaftenQuadratintegrierbarerAbbildungen4}
$\mfa-\mfb$--messbar und nach
\ref{ErwartungswertDer2NormEinesStandardnormalverteiltenProzesses} 
auch integrierbar, also ist es auch $f_Y$ als Vielfaches einer integrierbaren
Abbildung.\\
Seien $t \in T$, $v \in V$ und $\omega \in \Omega$.\\
Dann gilt:
\begin{align*}
&\abs{Y(t, v)(\omega)} \overset{\refclaim{BeschraenktheitGordonscherPaare}
{BeschraenktheitGordonscherPaare1}}\le 
 \opnorm{g(\omega)}{\R^m}{\R^N}\norm t _{2, m}\norm v_{2, N}
 \le\, ab\opnorm{g(\omega)}{\R^m}{\R^N}\\
  =&\, ab \opnorm{\Mat_{m, N}(\widetilde g (\omega))}{\R^m}{\R^N}
 = ab \norm{\widetilde g (\omega)}_{\Mat_{m, N}} \overset{(\ast\ast)}\le
 abd \norm{\widetilde g (\omega)}_{2, mN} = f_Y(\omega) \text.
\end{align*}
Somit ist auch (2) (ohne den Zusatz) gezeigt.\\
Seien $(E, \norm\cdot_E)$ ein normierter $\R$--Vektorraum, $x \in E^N$ und
$G_x$ der Gauß'sche Operator zum Tupel $(x,g)$.\\
Sei $V'':= \menge{\pfeil \vi N(x)}{\vi \in \kabg{E'}01}$.\\
Wie im Beweis des Zusatzes von (1) gesehen, ist $V''$ beschränkt und
nichtleer.\\
Es gilt für alle $s\in T$ und alle $\omega \in \Omega$:
\begin{align*}
\sup_{v \in V''}Y(s,v) &= \sup_{v\in V''}\sprod{g(\omega)s}{v}_N
=\sup_{\vi \in \kabg{E'}01}\sprod{g(\omega)s}{\pfeil\vi N (x)}_N\\
&\overset{\refclaim{EigenschaftenDesLinearkombinators}
{EigenschaftenDesLinearkombinators2}}=
\sup_{\vi \in \kabg{E'}01}\vi\left(\bprod{g(\omega)s}x_N \right)
\overset{\ref{NormUndDualeNorm}}= \norm{\bprod{g(\omega)s}x_N}_E\\
&= \norm{G_x(\omega)s}_E 
= \left(\norm\cdot_E\circ ev(s,\cdot)\circ G_x\right)(\omega)\text.
\end{align*}
Somit gilt:
\begin{align*}
\inf_{t \in T}\, \norm\cdot_E \circ ev(t,\cdot)\circ G_x
&= \inf_{t \in T}\sup_{v \in V''}Y(t, v)\text,\\
\sup_{t \in T}\, \norm\cdot_E \circ ev(t,\cdot)\circ G_x
&= \sup_{t \in T}\sup_{v \in V''}Y(t, v)\end{align*}
und damit sind diese beiden Abbildungen nach (2) integrierbar.
\end{proof}